\begin{table}[t]
\centering
\small
\setlength{\tabcolsep}{6pt}
\renewcommand{\arraystretch}{1.10}

\begin{tabular}{lccc}
\toprule
\textbf{Held-out Disease} & \textbf{Trained Head AUROC} & \textbf{Trained Head AUPR} & \textbf{Energy AUROC} \\
\midrule

BTC     & 0.782 & 0.855 & 0.591 \\
CRC     & 0.874 & 0.901 & 0.779 \\
GC      & 0.729 & 0.698 & 0.575 \\
HCC     & 0.881 & 0.842 & 0.558 \\
Normal  & 0.607 & 0.843 & 0.722 \\
PC      & 0.999 & 0.999 & 0.539 \\
\midrule
\textbf{Macro-average} & \textbf{0.812} & \textbf{0.856} & \textbf{0.627} \\

\bottomrule
\end{tabular}

\caption{
Novelty detection under the nested leave-one-disease-out protocol. For
each held-out disease, the trained novelty head is learned on a
proxy-novel disease drawn from the remaining training diseases and then
evaluated on the truly held-out disease, which is unseen during both
cascade and novelty-head training. The trained head substantially
outperforms the strongest training-free baseline on average, but
performance varies across diseases, showing that unseen-disease novelty is
learnable yet heterogeneous in this small-cohort setting.
}
\label{tab:novelty}

\end{table}
